\subsection{Frozen predictions as auxiliary information}
\label{sec:model-assisted}

A probability sample can correct errors in telemetry predictions of total labour. The established difference estimator combines the predictions with a weighted sum of observed residuals. Denis and Haddou give the fixed-prediction form and its design variance explicitly \cite[Section~2.3, equations~(2)--(4)]{denishaddou2026}. Prediction-powered inference also identifies its mean construction as a difference estimator and considers finite-population extensions \cite[Section~1.4 and Appendix~B of the accompanying preprint]{angelopoulos2023}.

Let $U=\{1,\ldots,M\}$ be a complete, fixed frame of person-intervals, with true labour $y_m=d_ma_m$ measured in person-hours and total $H=\sum_Uy_m$. For every frame unit, retain a prediction $f_m$ in person-hours. A raw SCH value requires a stated conversion before it can serve as $f_m$, so that prediction and observed labour have the same units. The entire prediction vector must be fixed with respect to the realised sample. Predictions trained on an independent earlier population can qualify, conditional on their frozen values and the subsequent sampling design. Accurate labels, a complete frame and correct inclusion probabilities remain substantive requirements.

\begin{proposition}[Fixed-prediction difference estimator]
\label{prop:difference-estimator}
For a probability sample $s\subseteq U$ with known first-order inclusion probabilities $\pi_m>0$, set $r_m=y_m-f_m$ and
\begin{equation}
 \widehat H_{\mathrm{diff}}
 =\sum_{m\in U}f_m+\sum_{m\in s}\frac{y_m-f_m}{\pi_m}.
 \label{eq:difference-estimator}
\end{equation}
If $y_m$ and $f_m$ are fixed under the design and every selected label is observed accurately, then $\E_s[\widehat H_{\mathrm{diff}}]=H$. This holds without independence between inclusion indicators or correct prediction of individual outcomes.
\end{proposition}

\begin{proof}
Write $I_m=\ind\{m\in s\}$. Linearity and $\E_s[I_m]=\pi_m$ give
\[
 \E_s[\widehat H_{\mathrm{diff}}]
 =\sum_Uf_m+\sum_U\frac{\E_s[I_m]}{\pi_m}(y_m-f_m)
 =\sum_Uy_m.
\]
\end{proof}

The residual correction makes the estimate design-unbiased even when the prediction total is biased. Taking $f_m=0$ recovers the Horvitz--Thompson (HT) estimator. Selection must include observed zeros as well as positive labour. A selected unanswered prompt is nonresponse, which must remain distinguishable from zero labour and from an unselected unit. The design-unbiasedness guarantee requires a complete frame, accurate allocations and known sampling probabilities. A provenance declaration alone cannot establish these properties. Independent label error with zero conditional mean may permit unbiasedness under both sampling and measurement error, but that requires an additional assumption beyond the fixed-label result proved here.

\subsection{Precision depends on residuals and the sampling design}
\label{sec:difference-variance}

Let $\pi_{mk}=\E_s[I_mI_k]$, including $\pi_{mm}=\pi_m$. The exact design variance is
\begin{equation}
 \Var_s(\widehat H_{\mathrm{diff}})
 =\sum_{m,k\in U}
 \frac{\pi_{mk}-\pi_m\pi_k}{\pi_m\pi_k}\,r_mr_k.
 \label{eq:difference-general-variance}
\end{equation}
This follows by expanding the variance of $\sum_U I_mr_m/\pi_m$. The covariances require joint inclusion probabilities. First-order probabilities alone cannot determine them. Equation~\eqref{eq:difference-general-variance} is a population identity whose unsampled residuals are generally unknown. Estimating its value requires the actual sampling design and the joint-inclusion conditions needed by the variance estimator.

For simple random sampling without replacement of exactly $n$ units from $M\geq2$, every size-$n$ subset has the same probability and $\pi_m=n/M$. Write
\[
 S_r^2=\frac{1}{M-1}\sum_U(r_m-\bar r_U)^2,
 \qquad
 S_y^2=\frac{1}{M-1}\sum_U(y_m-\bar y_U)^2.
\]
The inclusion covariances in \eqref{eq:difference-general-variance} give
\begin{equation}
 \Var_s(\widehat H_{\mathrm{diff}})
 =\frac{M^2}{n}\left(1-\frac nM\right)S_r^2,
 \qquad
 \Var_s(\widehat H_{\mathrm{HT}})
 =\frac{M^2}{n}\left(1-\frac nM\right)S_y^2.
 \label{eq:difference-srs-variance}
\end{equation}
For $2\leq n\leq M$, the unbiased sample-variance estimator is
\begin{equation}
 \widehat V_{\mathrm{diff}}
 =\frac{M^2}{n}\left(1-\frac nM\right)s_r^2,
 \qquad
 s_r^2=\frac{1}{n-1}\sum_{m\in s}(r_m-\bar r_s)^2.
 \label{eq:difference-srs-variance-estimator}
\end{equation}
Indeed, $s_r^2=[n(n-1)]^{-1}\sum_{m<k,\,m,k\in s}(r_m-r_k)^2$. Each pair has inclusion probability $n(n-1)/[M(M-1)]$, so $\E_s[s_r^2]=S_r^2$. A noncensus sample of size one does not supply this variance estimate. A complete census has zero sampling variance, including a single-unit population. Measurement error remains a separate question.

For $n<M$, the difference estimator improves on HT exactly when $S_r^2<S_y^2$. The gain in sampling precision therefore depends on residual variation under the design. Consider $M=4$, $n=2$, and $y=(0,2,4,6)$. The true total is $12$, and HT variance is $80/3$. Frozen predictions $f=(1,3,3,5)$ yield variance $16/3$, whereas $f=(6,4,2,0)$ yield $320/3$. Enumerating all six samples reproduces these variances and unbiased totals. It also verifies that \eqref{eq:difference-srs-variance-estimator} averages to each corresponding variance. The second predictor increases variance fourfold. These are finite mathematical checks, not evidence that a code-derived predictor improves precision in a workplace.

\subsection{Failure cases and implementation checks}
\label{sec:difference-contract}

Fitting predictions on the realised sample invalidates the fixed-vector proof. For $M=2$, $n=1$ and $y=(0,2)$, suppose a fitted rule memorises the selected outcome and predicts zero for the other unit. Its observed residual is always zero. The difference total is nevertheless $0$ or $2$, each with probability $1/2$, giving expectation $1$ rather than $H=2$. Generic cross-fitting is not a substitute for checking the conditioning and inclusion probabilities under the actual finite-population design.

The estimate can also be negative even though true labour is nonnegative. With $y=(0,0)$, $f=(10,0)$ and one uniformly selected unit, \eqref{eq:difference-estimator} is $-10$ or $10$ with equal probability. Clipping at zero changes its expectation from $0$ to $5$. Report the signed estimate and investigate an implausible result rather than silently changing the estimator. An unbiased labour denominator likewise does not make retained output divided by that denominator an unbiased productivity ratio.

The observatory's \texttt{estimate-labour} command requires a declaration that the frame is complete, stable unit identities, frozen predictions with their provenance, and positive inclusion probabilities. Each unit is recorded as unselected, observed or selected with nonresponse. It rejects nonresponse, missing freeze or frame attestations, duplicate identities, invalid probabilities and nonfinite arithmetic. A supplied general design supports point estimates only. Declared simple random sampling additionally requires the recorded selected count and probabilities to agree with $n$ and $M$. It supplies \eqref{eq:difference-srs-variance-estimator} for $n\geq2$ or zero sampling variance for a census. These consistency checks cannot establish that the declared sampling procedure was followed.

The implementation preserves signed estimates, observed zeros and unavailable uncertainty. It reports variance and standard error without a default confidence interval. An unbiased variance estimator alone establishes neither a normal approximation nor finite-sample interval coverage. Exhaustive synthetic tests cover unequal-probability sampling, both useful and harmful predictions, the sample-fitting counterexample, census and singleton cases, and invalid or missing observations. They verify the finite examples, input checks and reported estimates. Frame completeness, label accuracy and empirical calibration require evidence from the measurement process.
